Single-cell transcriptomic analysis requires both a partition of cells and an explanation of the genes associated with that partition. This report reproduces the correspondence-analysis and directional-clustering workflow of CAdir and evaluates its behavior on two mouse datasets. The original-paper case uses the Limb Muscle subset of Tabula Muris. The input contains 13,204 genes and 1,882 cells. Existing log-normalized expression values are retained, 6,000 highly variable genes are selected, and correspondence analysis provides a 30-dimensional representation. CAdir is initialized with eight directions, a cutoff angle of 55 degrees, and random seed 2358. The analysis returns seven cell clusters and 1,200 assigned genes. Against the original reference annotations, the adjusted Rand index is 0.669519 and normalized mutual information, calculated using the larger of the two entropies, is 0.741885.

The LM analysis implements the four panels of the published Figure 3 and five associated supplementary panels. These outputs include association plots, immune and endothelial direction comparisons, an actual split-and-merge graph, pairwise direction angles, a direction projection matrix, a reference-label Sankey diagram, and comparisons between principal and standard coordinates. The published author execution record agrees with the reproduced cell count, assigned-gene count, final cluster count, all reported cluster sizes and names, and the visible split-and-merge messages. The adjusted Rand index also agrees at the precision printed by the author. Independent checks cover saved coordinates, gene scores and rankings, clustering metrics, \seqsplit{overrepresentation} analysis, annotation assignment, and graph topology.

The extension experiment uses the Nestorowa hematopoietic stem and progenitor cell study. A common preprocessing procedure retains 1,656 cells and selects 6,000 highly variable genes. External evaluation uses the same fixed set of 1,641 cells carrying published broad immunophenotype labels. CAdir is compared with a principal-component-analysis, rank-based shared-nearest-neighbor graph, and weighted Louvain workflow. Both methods use 30-dimensional representations and the same five predefined seeds. All ten runs complete successfully. Median adjusted Rand indices are 0.291506 for CAdir and 0.269345 for the baseline. Median normalized mutual information values are 0.343414 and 0.356464, respectively. The metric ordering therefore depends on the criterion used. CAdir produces nine to eleven clusters, whereas the baseline produces seven in each run. This difference provides a concrete basis for examining how cluster granularity affects agreement with the reference partition.

The report connects these results to mathematical definitions and executable source code. It explains correspondence-analysis coordinates, directional assignment, total-least-squares direction updates, gene selection and ranking, and reference-based evaluation. The execution procedure records exact input identifiers, configuration, source revisions, random-number states, saved analysis objects, and resource measurements. Figures are generated from the saved objects, linking each displayed result to the corresponding numerical tables. Together, the original-paper case and the new-data comparison demonstrate an interpretable analysis workflow in which clustering, associated genes, geometric diagnostics, and quantitative evaluation can be inspected as connected parts of one experiment.
\par
